\part{系统总体方案设计}
\section{系统总体方案设计}
\subsection{系统架构体系}
基于传统气垫船操控复杂度高、状态感知能力有限及智能化程度不足的现状，本研究设计并实现了一套基于 CH32 的智能气垫船控制系统。该系统创新性地集成了多模态姿态融合控制与自主决策单元，通过引入多种环境与机体传感器，实时采集航行状态与周围环境数据。构建了包括姿态稳定、高度保持、路径跟踪在内的闭环控制系统，并结合故障诊断与应急策略模块，实现对推进系统与悬浮系统的智能管理与安全保护。此外，系统搭载了远程数传与图像传输模块，支持航行数据与实时画面的远程回传，并开发了跨平台地面站软件，为用户提供直观的可视化监控界面与便捷的指令下发功能。围绕上述功能设计，本章将对气垫船控制系统的整体架构进行详细阐述。
\subsubsection{需求分析} 
针对当前气垫船在功能完善性与监控能力方面的不足，本文在优化系统稳定性的同时，对系统功能进行了拓展，提出了以下改进方向及实现要点：
\begin{enumerate}[i)]
    \item 动力与运动控制：使用两个电机驱动风扇，让船体悬浮起来。使用四个电机驱动螺旋桨，控制船体前进、后退、平移和原地旋转。
    \item 自动姿态稳定：通过传感器实时监测船身的倾斜和转向角度。自动调节四个推进电机的力量，保持船体平稳，防止侧翻。
    \item 图形化显示与设置：在屏幕上直观显示船体姿态、速度、电量、电机状态等信息。可以直接在触屏上修改和调整各种运行参数。
    \item 故障诊断与报警：实时监测系统状态，当发生电机异常、传感器故障、通讯中断等问题时，屏幕会立即显示警告信息。在危险情况下，会发出报警并自动限制动力，保障安全。
    \end{enumerate} 
% 详细说明系统必须完成的动作，例如：通过无刷电机实现稳定气垫悬浮、基于视觉传感器实现自主赛道循迹、实时解算并控制船体姿态、通过交互菜单在线调整参数等。定义量化标准，例如：悬浮高度控制精度（如±5mm）、最大直线速度（如2m/s）、控制周期（如10ms）、连续运行时间等。

\subsubsection{系统总体框架}
以框图形式展示以CH32V307主控为核心，连接感知单元（摄像头、IMU）、执行单元（悬浮/推进无刷电机及其驱动模块）、交互单元（屏幕、按键）以及供电单元的物理架构和数据流向。
\subsection{控制系统概述}
\subsubsection{基本运动状态定义} 
定义船体的关键运动状态，如：静止悬停、直线加速、弯道循迹、紧急制动等，并分析各状态下执行器的协同工作模式。
\subsubsection{多电机协同控制框架}
阐述用于悬浮、推进及可能转向的无刷电机之间，如何根据不同的运动状态进行速度和力矩的分配与协调。
\subsubsection{视觉传感器数据处理流程} 器...”
概述摄像头图像采集、赛道特征识别（如边缘、中线）、以及路径偏差计算的完整算法流程。
\subsubsection{惯性测量单元数据融合}
概述IMU原始数据（加速度、角速度）的读取、滤波（可提及卡尔曼滤波）以及最终姿态角（俯仰、横滚）的解算流程。
\subsubsection{运行参数管理与存储} 
说明PID参数、控制阈值等关键运行参数的存储（如Flash）与通过四级菜单进行在线修改、保存和加载的机制。
\subsubsection{系统状态监控与显示}
描述如何将船体实时状态（如速度、姿态角、电池电压）在交互界面上进行可视化和更新。

\subsection{远程遥控器通信概述}
\subsubsection{网络基本架构}
\subsubsection{协议帧格式}
\subsubsection{通信协议}

\subsection{系统整体设计方案}
\subsubsection{设备功能层（硬件执行层）} % 对应“2.5.1 设备功能层”
归纳所有传感器、执行器、主控芯片等物理设备构成的基础功能层。
\subsubsection{控制传输层（算法决策层）} % 对应“2.5.2 控制传输层”
归纳FreeRTOS任务调度、通信及核心控制算法构成的决策与信息流转层。
\subsubsection{交互管理层（应用接口层）} % 对应“2.5.3 设备管理展示层”
归纳参数管理、状态显示等与人交互的应用层功能。

\subsection{本章小结}
总结本章提出的气垫船控制系统总体方案的核心内容，包括从需求到架构、从硬件到软件、从通信到算法的多层次设计，并指出该方案为后续第三、第四章的详细实现提供了明确的蓝图和依据。